% マクスウェル方程式 --- LuaLaTeX でコンパイル
\documentclass[lualatex, a4paper, 12pt]{jlreq}

% 数学パッケージ
\usepackage{amsmath}    % 高度な数式環境
\usepackage{amssymb}    % 数学記号
\usepackage{bm}         % 太字ベクトル \bm{...}

% ページ設定（オプション）
\usepackage{geometry}
\geometry{top=25mm, bottom=25mm, left=25mm, right=25mm}

\title{マクスウェル方程式}
\author{電磁気学}
\date{\today}

\begin{document}
\maketitle

\section{概要}

電磁気学はマクスウェル方程式と呼ばれる4つの方程式によって
完全に記述される．
電場 $\bm{E}$，磁場 $\bm{B}$，電荷密度 $\rho$，
電流密度 $\bm{J}$ を用いて表される．

\section{積分形}

\begin{align}
  \oint_S \bm{E} \cdot d\bm{A}
    &= \frac{Q}{\varepsilon_0}
    \tag{ガウスの法則}
    \label{eq:gauss_e}
    \\[8pt]
  \oint_S \bm{B} \cdot d\bm{A}
    &= 0
    \tag{磁気のガウスの法則}
    \label{eq:gauss_b}
    \\[8pt]
  \oint_C \bm{E} \cdot d\bm{l}
    &= -\frac{d}{dt} \int_S \bm{B} \cdot d\bm{A}
    \tag{ファラデーの法則}
    \label{eq:faraday}
    \\[8pt]
  \oint_C \bm{B} \cdot d\bm{l}
    &= \mu_0 \int_S \bm{J} \cdot d\bm{A}
       + \mu_0 \varepsilon_0 \frac{d}{dt}
         \int_S \bm{E} \cdot d\bm{A}
    \tag{アンペール--マクスウェルの法則}
    \label{eq:ampere}
\end{align}

\section{微分形}

ガウスの定理とストークスの定理を用いると，
積分形から微分形が導かれる．

\begin{align}
  \nabla \cdot \bm{E}
    &= \frac{\rho}{\varepsilon_0}
    \tag{ガウスの法則}
    \label{eq:gauss_e_diff}
    \\[8pt]
  \nabla \cdot \bm{B}
    &= 0
    \tag{磁気のガウスの法則}
    \label{eq:gauss_b_diff}
    \\[8pt]
  \nabla \times \bm{E}
    &= -\frac{\partial \bm{B}}{\partial t}
    \tag{ファラデーの法則}
    \label{eq:faraday_diff}
    \\[8pt]
  \nabla \times \bm{B}
    &= \mu_0 \bm{J}
       + \mu_0 \varepsilon_0 \frac{\partial \bm{E}}{\partial t}
    \tag{アンペール--マクスウェルの法則}
    \label{eq:ampere_diff}
\end{align}

\section{物理定数}

\begin{table}[h]
  \centering
  \caption{電磁気学の基本定数}
  \begin{tabular}{cll}
    \hline
    記号 & 名称 & 値 \\
    \hline
    $\varepsilon_0$ & 真空の誘電率 &
      $8.854 \times 10^{-12}\ \mathrm{F/m}$ \\
    $\mu_0$ & 真空の透磁率 &
      $4\pi \times 10^{-7}\ \mathrm{H/m}$ \\
    $c$ & 真空中の光速 &
      $c = 1/\sqrt{\mu_0 \varepsilon_0}
       \approx 3.0 \times 10^8\ \mathrm{m/s}$ \\
    \hline
  \end{tabular}
  \label{tab:constants}
\end{table}

\section{電磁波}

マクスウェル方程式から，電場と磁場が波として伝播すること
（電磁波）が導かれる．電場の波動方程式：

\begin{equation}
  \nabla^2 \bm{E}
  - \mu_0 \varepsilon_0 \frac{\partial^2 \bm{E}}{\partial t^2}
  = \bm{0}
\end{equation}

この波の速度は $c = 1/\sqrt{\mu_0\varepsilon_0}$ であり，
光速と一致する．これがマクスウェルによる
「光は電磁波である」という発見である．

\end{document}
